\section{CRT decoupling and modular non-locality}
\label{sec:crt_decoupling}

Sections~\ref{sec:field_connection} and~\ref{sec:holonomy_primitive}
have established the status of the holonomy angle $\theta_p$ as a
primitive invariant and that of the classical objects $A_\mu,
F_{\mu\nu}, g_{\mu\nu}$ as derivatives of the scalar potential
$S$. To connect this status to the experimental question of
non-locality, it remains to examine the correlations between
distinct modular channels --- the only ``inter-object''
correlations that PT allows to carry a non-trivial informational
significance. This section states, without detailed proof, the
two structural theorems proved in the technical companion
paper~\cite{PT_CRT_DECOUPLING}, together with their empirical
corollary.

\subsection{Theorem 4.1: Geometric Decoupling Theorem (empirical form)}
\label{ssec:thm_4_1_en}

\begin{theorem}[Geometric Decoupling Theorem, empirical form]
\label{thm:invariance_geom_en}
\footnote{Formally verified in Lean~4 + Mathlib as
\texttt{PT.CrtDecoupling.Empirical.empirical\_invariance}
(hypothesis form),
\texttt{SpectralReduction.empirical\_invariance\_closed}
(finite-state closed form),
\texttt{Phase3.empirical\_invariance\_infinite}
(infinite-state form via Birkhoff--Hopf), and
\texttt{Phase4.sieve\_empirical\_invariance}
(sieve-specific closure);
see \cite[\S~8.5]{PT_CRT_DECOUPLING}.}
Let $p, q \in \{3, 5, 7\}$ be distinct active primes and
$\Phi_p \in \mathcal{A}_p$, $\Phi_q \in \mathcal{A}_q$ bounded
observables on the corresponding modular $\sigma$-algebras. The
function
\begin{equation}
  \mathcal{I}_{p,q}(\mu)
  \;:=\;
  I_{\pi_m^{\mathrm{emp}}}\!\bigl(\Phi_p\,;\,\Phi_q\,\big|\,G(\mu)\bigr)
\end{equation}
is constant on $\mu \in [\mu_c, \infty)$.
\end{theorem}

The full proof is given as Theorem~6.1 of the technical
companion paper~\cite{PT_CRT_DECOUPLING}. It proceeds in four
steps: (i) the empirical measure $\pi_m^{\rm emp}$ of the residue
trajectory is a functional of the sole field $\{g_n\}$, hence
independent of $\mu$; (ii) the geometry $G(\mu)$ is measurable
with respect to the $\sigma$-algebra generated by the spectrum of
the transfer matrix $T_m$, which is itself contained in the
shift-invariant $\sigma$-algebra; (iii) the sieve dynamics is
ergodic by Birkhoff--Hopf applied to the spectral gap T4, hence
the invariant $\sigma$-algebra is trivial modulo
$\pi_m^{\rm emp}$, hence $G(\mu)$ is almost surely constant;
(iv) conditioning on an almost surely constant random variable
does not change mutual information. The conclusion follows.

The conceptual point is that the emergent metric $G(\mu)$ ---
including across the Hartle--Hawking transition at $\mu_c \approx
6.97$ and up to the canonical fixed point $\mustar = 15$ --- never
modifies the joint informational content of the active modular
channels. Inter-channel correlations are \emph{intrinsically}
pre-geometric: they exist before the construction of the metric
and are indifferent to it.

\subsection{Theorem 4.2: Geometric Decoupling Theorem (Ruelle form)}
\label{ssec:thm_4_2_en}

\begin{theorem}[Geometric Decoupling Theorem, Ruelle form]
\label{thm:ruelle_zero_en}
\footnote{Formally verified in Lean~4 + Mathlib as
\texttt{PT.CrtDecoupling.Main.geometric\_decoupling\_ruelle};
see \cite[\S~8.5]{PT_CRT_DECOUPLING}.}
Under the idealised measure $\pi_m^{\mathrm{Ruelle}}$, we have
$I(\Phi_p\,;\,\Phi_q) = 0$ exactly.
\end{theorem}

The proof is immediate from the CRT tensor factorisation of the
transfer $T_{pq} = T_p \otimes T_q$ (Theorem~3.1
of~\cite{PT_CRT_DECOUPLING}) and the corresponding factorisation
of the unique left Perron eigenvector
$\pi_{pq}^{\rm Ruelle} = \pi_p^{\rm Ruelle} \otimes \pi_q^{\rm Ruelle}$
(Lemma~3.2 of the same paper). This factorisation is an
independence condition, hence the Kullback--Leibler divergence
between the joint law and the product of marginals is zero; by
the data-processing inequality, the mutual information of any
pair of bounded observables inherits the vanishing.

\emph{Numerical validation.} The factorisation
$\pi_{15}^{\rm Ruelle} = \pi_3^{\rm Ruelle} \otimes \pi_5^{\rm Ruelle}$
is verified at the residual level
\begin{equation}
  \|\pi_{15}^{\rm direct} - \pi_3 \otimes \pi_5\|_2
  \;=\; 5.88 \times 10^{-53}
\end{equation}
at \texttt{mpmath} 50 digits of precision (test T2 PASS of
\texttt{scripts/test\_crt\_decoupling.py}, see
\cite{PT_CRT_DECOUPLING}).

\subsection{Corollary 4.3: Decoupling Closed Form}
\label{ssec:cor_4_3_en}

\begin{corollary}[Decoupling Closed Form]
\label{cor:closed_form_en}
\footnote{Not formalisable as a Lean theorem in the present form
because the depth parameter $k_{p,q}^{\rm eff}$ is fitted
empirically; only the prediction at a fixed $k$ admits a
Lean-verifiable computational statement;
see \cite[\S~8.5]{PT_CRT_DECOUPLING}.}
The constant $\mathcal{I}_{p,q}$ admits a closed-form expression
derived from the holonomy angles:
\begin{equation}
  \mathcal{I}_{p,q}
  \;=\;
  \frac{1}{2 \ln 2}\Bigl[
    \cos^{2 k_{p,q}}(\theta_3, \mu^\ast)
    + \sin^{4 k_{p,q}}(\theta_2, p_1)
  \Bigr].
  \label{eq:closed_form_en}
\end{equation}
Numerical validations (cf.\ \cite{PT_CRT_DECOUPLING}, Table~B):
$\mathcal{I}_{3,5} = 0.132$ bits (target $0.138$, error $4.5\%$);
$\mathcal{I}_{3,7} = 0.306$ bits (target $0.283$, error $8.1\%$).
\end{corollary}

The depth parameter $k_{p,q}$ is fitted empirically at
$k_{3,5}^{\rm eff} = 7$ and $k_{3,7}^{\rm eff} = 4$; a
combinatorial derivation from the antidiagonal structure of $T_3$
remains open (cf.\ Remark~7.1 of \cite{PT_CRT_DECOUPLING}). This
uncertainty affects only the \emph{quantitative value} of
$\mathcal{I}_{p,q}$; the structural result
(Theorem~\ref{thm:invariance_geom_en}) is independent of the
choice of $k_{p,q}$.

\subsection{Spatial non-locality dissolved into modular non-locality}
\label{ssec:dissolution_en}

The three preceding statements allow one to reformulate, in PT
terms, the status of the apparent ``non-locality'' of quantum
correlations. One can summarise the situation in three propositions
strictly parallel to the three positions of the foundational
debate (\S~\ref{sec:debate}), but with a different status:

\paragraph{(i) Inter-channel correlations exist.}
$\mathcal{I}_{3,5}$ and $\mathcal{I}_{3,7}$ are non-zero and of the
magnitude predicted by Corollary~\ref{cor:closed_form_en}. PT does
not deny the existence of EPR/Bell correlations; on the contrary,
it provides a quantitative prediction of them.

\paragraph{(ii) These correlations are of arithmetic origin.}
The closed form of Corollary~\ref{cor:closed_form_en} involves
only the holonomy angles $\theta_3, \theta_2$ evaluated at the
canonical scales $\mustar = 15$ and $p_1 = 3$. No spatial
integral, no metric, no distance between points of space enters
the expression. The observed non-locality is not a non-locality
in space: it is an intra-CRT correlation between different
modular channels.

\paragraph{(iii) These correlations are invariant under $G(\mu)$.}
Theorem~\ref{thm:invariance_geom_en} forbids the emergent
geometry from \emph{coupling} the modular channels. Whatever the
position in the geometric family $\mu \in [\mu_c, \infty)$ ---
including at any distance between detectors in a Bell
configuration --- the mutual information between distinct
channels remains the same.

These three propositions together dissolve spatial non-locality
in the strong sense of position B: quantum correlations do not
cross space because they are not in space. They are in the CRT
structure, which manifests in the emergent geometry $G$ by
projection without crossing it or using it as a causal carrier.
The apparent paradox of information ``faster than $c$'' in Bell
tests is an artefact of the spatial reading of a correlation that,
in the PT substrate, is not spatial. We return to the experimental
consequences of this dissolution in
\S~\ref{ssec:bell_consequence_en}.
